% !TEX root = ../thesis.tex

\chapter{Methodology} \label{chp:method}

\section{Robot Platform: B2 Quadruped + Z1 Arm}

The hardware platform used in this project is the Unitree B2 quadruped body
paired with the Unitree Z1 robotic arm, forming the \textbf{B2+Z1} whole-body
system simulated in NVIDIA Isaac Gym.

\subsection{Unitree B2 Quadruped}

The B2 is a large-scale, high-payload quadruped with the following characteristics
relevant to control design:
\begin{itemize}
  \item \textbf{Body mass:} approximately 52\,kg (compared to Go1 at $\approx$12\,kg
        and Spot at $\approx$32\,kg).
  \item \textbf{Degrees of freedom:} 12 revolute joints (3 per leg: hip abduction/adduction,
        hip flexion/extension, knee flexion/extension).
  \item \textbf{PD joint control:} each joint is driven by a position controller
        $\tau = K_p(\theta^* - \theta) - K_d\dot{\theta}$, where $\theta^*$ is the
        policy-output target position.  The heavier body mass requires substantially
        higher gains than the Go1: we use $K_p = 200\,\text{N\,m/rad}$ and
        $K_d = 20\,\text{N\,m\,s/rad}$ for all leg joints.
\end{itemize}

\subsection{Unitree Z1 Arm}

The Z1 is a 6-DoF lightweight arm (payload $\approx$2\,kg) mounted on the front
of the B2 torso.
The six controlled joints are, in order from base to tip:
\texttt{z1\_waist}, \texttt{z1\_shoulder}, \texttt{z1\_elbow},
\texttt{z1\_wrist\_angle}, \texttt{z1\_forearm\_roll},
\texttt{z1\_wrist\_rotate}.
Each joint is also controlled by a PD servo; the gains are set per-joint
(e.g.\ $K_p^{\text{shoulder}}=60$\,N\,m/rad) to reflect the varying
load-bearing requirements along the kinematic chain.
The end-effector (index 23 in the URDF link tree) is the controlled point for
pose tracking.

\section{Framework Adaptation from Go1 to B2+Z1}
\label{sec:framework_adaptation}

The RoboDuet codebase was originally designed and validated on the
Unitree Go1 + Z1 system~\cite{pan2024roboduet}.
Adapting it to the B2+Z1 required the following changes.

\subsection{URDF and Asset Configuration}

The simulation asset file was changed from
\texttt{go1/urdf/go1.urdf} to
\texttt{b2z1/urdf/b2z1.urdf}.
The B2+Z1 combined URDF provides the complete kinematic and inertial
description of both the quadruped body and the mounted arm,
including mesh geometries for collision and visual rendering.
The Isaac Gym environment was updated to load this asset, and the
relevant link indices (feet, base, end-effector at link 23) were
re-mapped to match the B2+Z1 joint ordering.

\subsection{PD Gain Correction}
\label{subsec:pdgain}

The most critical adaptation was correcting the PD servo gains.
The Go1 framework uses $K_p = 35$\,N\,m/rad and $K_d = 1$\,N\,m\,s/rad,
which are appropriate for the lightweight Go1 body.
Applying these gains directly to the B2---which is approximately four times
heavier---produced oscillatory, unstable behaviour in early experiments:
the robot's legs would undershoot commanded joint positions and the body
would collapse.
Increasing the gains to $K_p = 200$\,N\,m/rad and $K_d = 20$\,N\,m\,s/rad
(a $5.7\times$ increase in stiffness and $20\times$ increase in damping)
resolved this instability and produced stable, trotting locomotion within
the same number of training iterations as the Go1 baseline.

The corrected gains are applied consistently across both the locomotion
training stage and the manipulation training stage:
\begin{verbatim}
  Cfg.dog.control.stiffness_leg["joint"] = 200.0
  Cfg.dog.control.damping_leg["joint"]   = 20.0
  Cfg.control.stiffness["joint"]         = 200.0
  Cfg.control.damping["joint"]           = 20.0
\end{verbatim}

\subsection{Configuration Pipeline}

The configuration follows a three-step chain inherited from the
original framework:
\begin{enumerate}
  \item \texttt{config\_go1(Cfg)} --- sets the base locomotion defaults
        (reward scales, network sizes, domain randomisation ranges).
  \item \texttt{config\_wtw(Cfg)} --- overlays the Walk-These-Ways
        distributional command space~\cite{margolis2022walk}:
        gait frequency $[2,4]$\,Hz, phase offset, footswing height,
        stance width/length, body pitch/roll command ranges.
  \item \texttt{config\_asset(Cfg)} --- sets the B2+Z1 asset file path
        and link/joint index mappings.
\end{enumerate}
Per-experiment overrides (PD gains, episode length, observation size,
environment count) are applied after these three calls.

\section{Three-Stage Training Pipeline}
\label{sec:training_pipeline}

Training is organised into three sequential stages.
Each stage uses PPO~\cite{schulman2017proximal} with separate
actor-critic networks for the dog policy and the arm policy.
2048 parallel Isaac Gym environments are used throughout.
The simulator runs at 200\,Hz; policy inference and action
updates occur at 50\,Hz (every 4 simulation steps), so one
policy step corresponds to 20\,ms of simulated time.

\subsection{Stage 1: Locomotion Pre-training}

\textbf{Goal:} train the dog policy to produce stable, command-following
trotting gait on flat terrain.

\textbf{Configuration:}
\begin{itemize}
  \item Arm is \emph{fixed} in a resting configuration
        (\texttt{keep\_arm\_fixed = True}).
  \item Locomotion commands: forward velocity $v_x \in [-1.0, 1.0]$\,m/s,
        lateral velocity $v_y \in [-0.6, 0.6]$\,m/s,
        yaw rate $\omega_z \in [-1.0, 1.0]$\,rad/s.
  \item Gait parameters (frequency, phase, offset, bound, duration,
        foot-swing height, stance dimensions) are randomly sampled
        per episode from the distributional command curriculum.
  \item Duration: 50,000 PPO iterations, 2048 environments.
\end{itemize}

\textbf{Outcome:} the dog policy converges to a stable trotting gait;
the reward stabilises and the robot does not fall over across long rollouts.
The saved checkpoint \texttt{ac\_weights\_last\_dog.pt} is used as the
starting point for Stage 3.

\subsection{Stage 2: Arm Manipulation Pre-training (Stand-Still)}

\textbf{Goal:} train the arm policy to track random 6-DoF end-effector
pose commands while the robot body is \emph{stationary}.

\textbf{Configuration:}
\begin{itemize}
  \item Locomotion velocity commands are set to zero
        ($v_x = v_y = \omega_z = 0$), forcing the body to stand still.
  \item Arm is \emph{enabled} from the first iteration
        (\texttt{keep\_arm\_fixed = False},
        \texttt{pretrained\_to\_hybrid\_start = 0}).
  \item The arm manipulation reward weight is boosted to $5.0$
        (from the default $1.0$) to accelerate early arm learning.
  \item The body-height terminal condition (collapse height $0.28$\,m)
        is enabled to prevent the dog policy from losing balance during
        arm motion.
  \item Duration: 30,000 PPO iterations.
\end{itemize}

\textbf{Outcome:} the arm policy learns to reach commanded positions and
match commanded orientations, while the dog policy learns to compensate
for the arm's disturbance forces on the body.
The saved checkpoint \texttt{ac\_weights\_last\_arm.pt} is used as the
starting point for Stage 3.

\subsection{Stage 3: Hybrid Fine-tuning (Walk + Grasp)}

\textbf{Goal:} fine-tune both policies simultaneously so that the robot
can walk while executing arm manipulation.

\textbf{Configuration:}
\begin{itemize}
  \item Both the dog checkpoint (from Stage~1) and the arm checkpoint
        (from Stage~2) are \emph{hot-started} via:
        \texttt{DogRunnerArgs.resume = True},
        \texttt{ArmRunnerArgs.resume = True}.
  \item Locomotion commands are re-enabled with the same ranges as Stage~1.
  \item The arm manipulation reward weight is set to $3.0$
        (reduced from $5.0$ to balance locomotion and manipulation objectives).
  \item Locomotion tracking rewards are re-enabled at $70\%$ ($v_x$)
        and $50\%$ ($\omega_z$) of the Stage-1 scale, allowing the locomotion
        objective to share gradient signal with the manipulation objective.
  \item Duration: 30,000 PPO iterations.
\end{itemize}

\textbf{Outcome:} the robot is able to walk at commanded velocities while
simultaneously tracking end-effector pose targets with the Z1 arm.
The cooperative communication channel (body orientation command from arm
to dog) is active throughout this stage, enabling the body to lean as needed.

\section{Policy Architecture and Observation--Action Spaces}

The dual-policy architecture follows RoboDuet~\cite{pan2024roboduet} exactly.
Both policies are multi-layer perceptron (MLP) actor-critics with a
proprioceptive history encoder.

\subsection{Dog (Locomotion) Policy}

\paragraph{Observations.}
The dog policy receives a stacked history of the last 30 proprioceptive
observations (each of length 56), giving a total input dimension of
$56 \times 30 = 1680$.
Each single-step observation includes:
\begin{itemize}
  \item Body angular velocity: $[\omega_x, \omega_y, \omega_z]$ (3 values)
  \item Projected gravity vector: $[g_x, g_y, g_z]$ (3 values)
  \item Locomotion command: $[v_x^*, v_y^*, \omega_z^*,
        \phi_{\text{freq}}, \phi_{\text{phase}}, \phi_{\text{offset}},
        \phi_{\text{bound}}, \phi_{\text{dur}}, h_{\text{swing}},
        \phi^*, \rho^*, w_{\text{stance}}, l_{\text{stance}}]$ (13 values)
  \item Clock inputs encoding gait phase: $[\sin(\phi_1), \cos(\phi_1),
        \sin(\phi_2), \cos(\phi_2)]$ (4 values)
  \item Leg joint positions relative to default: $\mathbf{q}^{\text{leg}} \in \mathbb{R}^{12}$
  \item Leg joint velocities: $\dot{\mathbf{q}}^{\text{leg}} \in \mathbb{R}^{12}$
  \item Previous leg joint action: $\mathbf{a}^{\text{leg}}_{t-1} \in \mathbb{R}^{12}$
\end{itemize}

\textit{Cooperative input from the arm policy:}
In Stage 3, the body pitch and roll commands $[\phi^*, \rho^*]$ within the
command vector are \emph{replaced} by the values output by the arm policy.
This is the key communication channel: the arm policy can request the body
to tilt in any direction to improve end-effector reachability.

\paragraph{Actions.}
The dog policy outputs 12 leg joint position targets
$\mathbf{a}^{\text{leg}} \in \mathbb{R}^{12}$, which are converted to
joint torques by the PD controller.

\subsection{Arm Policy}

\paragraph{Observations.}
The arm policy receives a stacked history of the last 30 arm observations
(each of length 20), giving a total input dimension of $20 \times 30 = 600$.
Each single-step observation includes:
\begin{itemize}
  \item Arm joint positions: $\mathbf{q}^{\text{arm}} \in \mathbb{R}^{6}$
        (z1\_waist through z1\_wrist\_rotate)
  \item Arm joint velocities: $\dot{\mathbf{q}}^{\text{arm}} \in \mathbb{R}^{6}$
  \item End-effector pose command $(l, p, y, r_{ee}, p_{ee}, y_{ee})$:
        target reach length, polar angle, azimuthal angle, and end-effector
        roll/pitch/yaw (6 values)
  \item Current body pitch and roll from the dog state (2 values)
\end{itemize}

\paragraph{Actions.}
The arm policy outputs two sets of values:
\begin{enumerate}
  \item \textbf{Arm joint position targets:}
        $\mathbf{a}^{\text{arm}} \in \mathbb{R}^{6}$ --- the desired joint
        positions for the six Z1 joints, converted to torques by the arm PD
        controller.
  \item \textbf{Body orientation command:}
        $[{\phi}^*_{\text{arm}}, {\rho}^*_{\text{arm}}] \in \mathbb{R}^{2}$
        (body pitch and roll) --- injected into the dog policy's command
        vector in place of the user-specified orientation command.
\end{enumerate}

Total action dimension for the combined system: $12 + 6 = 18$.

\section{Reward Design}
\label{sec:rewards}

The overall reward at each step is the sum of dog rewards $r^{\text{dog}}$
and arm rewards $r^{\text{arm}}$, each computed separately but sharing
the same simulation state.

\subsection{Dog Locomotion Rewards}

\begin{itemize}
  \item \textbf{Velocity tracking} ($r_{\text{lin}}$):
        $r = \exp(-\|\mathbf{v}_{xy} - \mathbf{v}^*_{xy}\|^2 / \sigma_v)$,
        encouraging the body to match commanded forward and lateral velocity.
  \item \textbf{Yaw tracking} ($r_{\text{ang}}$):
        $r = \exp(-(\omega_z - \omega_z^*)^2 / \sigma_\omega)$,
        penalising heading deviation.
  \item \textbf{Gait contact shaping} ($r_{\text{contact}}$):
        shaped foot contact force and foot velocity rewards that encourage
        the commanded gait pattern, following Margolis and
        Agrawal~\cite{margolis2022walk}.
  \item \textbf{Foot clearance} ($r_{\text{clear}}$, weight $-30.0$):
        penalises deviation of foot height during swing from the commanded
        foot-swing height.
  \item \textbf{Body orientation control} ($r_{\text{ori}}$, weight $-10.0$):
        penalises the difference between actual body pitch/roll and the
        body orientation command received from the arm policy.
  \item \textbf{Raibert heuristic} ($r_{\text{rait}}$, weight $-10.0$):
        encourages symmetric foot placement to improve dynamic stability.
  \item \textbf{Smoothness penalties} ($r_{\text{smooth1}}, r_{\text{smooth2}}$,
        weights $-0.1$ each): first- and second-order differences in
        leg joint actions.
  \item \textbf{Energy penalty} ($r_{\text{energy}}$, weight $-4 \times 10^{-5}$):
        penalises leg joint torque $\times$ velocity.
  \item \textbf{Collision penalty} ($r_{\text{col}}$, weight $-5.0$):
        penalises undesired body--ground contacts (non-foot links).
  \item \textbf{Velocity z penalty} ($r_{v_z}$, weight $-0.02$):
        penalises vertical body velocity.
  \item \textbf{Angular velocity xy penalty} ($r_{\omega_{xy}}$, weight $-0.001$):
        penalises body pitch and roll rates.
  \item \textbf{Terminal condition:} episode terminates (and receives
        a large negative reward) if the body height falls below $0.28$\,m.
\end{itemize}

\subsection{Arm Manipulation Rewards}

\begin{itemize}
  \item \textbf{End-effector pose tracking} ($r_{\text{manip}}$, weight $5.0$
        in Stage~2, $3.0$ in Stage~3):
        \[
          r_{\text{manip}} = w_{\text{lpy}} \cdot
          \exp\!\left(-\|[l,p,y]_{\text{ee}} - [l,p,y]^*\|^2\right)
          + w_{\text{rpy}} \cdot
          \exp\!\left(-\|[\text{roll,pitch,yaw}]_{\text{ee}}
          - [\text{roll,pitch,yaw}]^*\|^2\right),
        \]
        where $w_{\text{lpy}} = 3$ and $w_{\text{rpy}} = 1$.
        The position component receives triple the weight of the orientation
        component to prioritise accurate reach before orientation fine-tuning.
  \item \textbf{Arm smoothness} ($r_{\text{arm-smooth1}}, r_{\text{arm-smooth2}}$):
        first- and second-order differences in arm joint actions, weighted
        at $5\times$ the corresponding leg smoothness scales.
  \item \textbf{Arm velocity penalty} ($r_{\text{arm-vel}}$):
        penalises arm joint velocities, weighted $10\times$ the leg scale.
  \item \textbf{Arm acceleration penalty} ($r_{\text{arm-acc}}$):
        penalises arm joint accelerations, weighted $10\times$ the leg scale.
  \item \textbf{Arm action rate penalty} ($r_{\text{arm-rate}}$):
        penalises large changes in arm actions between steps,
        weighted $10\times$ the leg scale.
  \item \textbf{Arm energy penalty} ($r_{\text{arm-energy}}$, weight $-4\times 10^{-5}$):
        penalises arm torque $\times$ velocity.
\end{itemize}

The arm's heavier smoothness and velocity penalties (relative to the
leg penalties) reflect the requirement that the arm move in a smooth,
controlled manner to avoid destabilising the body via inertial reactions.

\section{Domain Randomisation}

To improve robustness and facilitate potential future sim-to-real transfer,
the following parameters are randomised at the start of each episode:
\begin{itemize}
  \item \textbf{Ground friction:} static and dynamic friction coefficient
        uniformly sampled from $[0.1, 3.0]$.
  \item \textbf{Restitution:} sampled from $[0.0, 0.4]$.
  \item \textbf{Added payload mass:} random mass added to the base link,
        sampled from $[-2.0, +2.0]$\,kg.
  \item \textbf{Gravity disturbance:} a random gravity offset sampled from
        $[-1.0, +1.0]$\,m/s$^2$ is applied every $8$\,s
        for $0.99$\,s each pulse.
  \item \textbf{Motor strength:} each joint's output torque is scaled by a
        factor sampled from $[0.9, 1.1]$.
  \item \textbf{Motor offset:} a fixed bias sampled from $[-0.02, +0.02]$\,rad
        is added to each joint's measured position.
  \item \textbf{Action latency:} a 6-step ($120$\,ms) observation lag is
        applied to simulate communication and computation delays.
\end{itemize}

\section{Training Setup and Hyperparameters}

\begin{table}[h]
\centering
\caption{Training hyperparameters for the B2+Z1 system.}
\label{tab:hparams}
\begin{tabular}{ll}
\hline
\textbf{Parameter} & \textbf{Value} \\
\hline
Simulator & NVIDIA Isaac Gym \\
Parallel environments & 2048 \\
Simulation frequency & 200\,Hz \\
Control frequency & 50\,Hz \\
RL algorithm & PPO~\cite{schulman2017proximal} \\
Discount factor $\gamma$ & 0.99 \\
GAE $\lambda$ & 0.95 \\
Clip ratio $\epsilon$ & 0.2 \\
Learning rate (dog) & $3 \times 10^{-4}$ (with sigmoid schedule) \\
Learning rate (arm) & $3 \times 10^{-4}$ (with sigmoid schedule) \\
Mini-batch size & 512 \\
PPO epochs per update & 5 \\
Episode length & 20\,s \\
Stage 1 iterations & 50,000 \\
Stage 2 iterations & 30,000 \\
Stage 3 iterations & 30,000 \\
Leg PD gains & $K_p = 200$, $K_d = 20$ \\
Arm PD gains (shoulder/elbow) & $K_p = 60$, $K_d = 6$ \\
Arm PD gains (waist) & $K_p = 40$, $K_d = 2$ \\
Arm PD gains (wrist joints) & $K_p = 20$, $K_d = 1$ \\
\hline
\end{tabular}
\end{table}

\noindent The training was performed on a server equipped with a single NVIDIA
RTX~4090 GPU.
Model checkpoints were saved every 500 iterations; the final model used for
evaluation corresponds to the last saved checkpoint at the end of each stage.
Weights \& Biases (WandB) was used for logging training curves and reward
component breakdowns throughout all three stages.
